\begin{figure}[htbp]
\centering
\renewcommand{\arraystretch}{2.2}
\begin{tabular}{|>{\centering\arraybackslash}m{2.8cm}%
                |>{\centering\arraybackslash}m{3.2cm}%
                |>{\centering\arraybackslash}m{3.2cm}%
                |>{\centering\arraybackslash}m{3.2cm}|}
\hline
\rowcolor{gray!22}
\textbf{Bedingung}
  & \textbf{T1: Bugfixes}
  & \textbf{T2: Refactorings}
  & \textbf{T3: Features} \\
\hline
\cellcolor{gray!12}{\small\textbf{B1\newline Reaktive\newline Exploration}}
  & \cellcolor{red!5}{\small Baseline}
  & \cellcolor{red!5}{\small Baseline}
  & \cellcolor{red!5}{\small Baseline} \\
\hline
\cellcolor{gray!12}{\small\textbf{B2\newline Statisch\newline (CLAUDE.md)}}
  & \cellcolor{yellow!18}{\small HF1}
  & \cellcolor{yellow!18}{\small HF1}
  & \cellcolor{yellow!18}{\small HF1} \\
\hline
\cellcolor{gray!12}{\small\textbf{B3\newline Living\newline Knowledge Graph}}
  & \cellcolor{green!10}{\small HF1 $\cdot$ HF2}
  & \cellcolor{green!10}{\small HF1 $\cdot$ HF2}
  & \cellcolor{green!16}{\small HF1 $\cdot$ HF2 $\cdot$ \textbf{HF3}} \\
\hline
\end{tabular}

\smallskip
{\footnotesize\textit{Metriken je Zelle: Token-Verbrauch, Iterationsanzahl, Erstkorrektheit,
Kontextpräzision. HF3 (Retrieval-Transparenz) wird über alle B3-Zellen erhoben; T3 ist der
primäre Messkontext.}}

\caption{Evaluationsdesign: drei Versuchsbedingungen (B1 bis B3) gekreuzt mit drei
Aufgabentypen (T1 bis T3). Jede Zelle steht für einen Messdurchlauf.}
\label{fig:evaluation}
\end{figure}
